% =====================================================================
%  2bat-ud3-2.tex  —  SESSIÓ 2: El sostre d'un servei (descoberta)
% =====================================================================
\horatitol{Sessió 2 — El sostre d'un servei}%
{Un cas real: descriure cap a quin valor s'acosta l'ocupació d'un servei que creix però mai supera la seva capacitat.}

\subsection{Idea clau}

Avui treballem en grup, a les pissarres, com a investigadors. Us repartim la
\textbf{gràfica muda} de l'ocupació d'un servei (per exemple, les places ocupades
d'un centre al llarg dels mesos): creix de pressa al principi i després cada cop més
a poc a poc, \textbf{sense superar mai la capacitat màxima}. El repte: descriure cap a
quin valor s'acosta, sense que us doni cap nom matemàtic. L'heu de trobar vosaltres.

\begin{idea}
\textbf{El cas.} Un centre té $200$ places. L'ocupació al llarg dels mesos creix, però
mai no pot passar de $200$. La pregunta del dia:
\begin{itemize}[leftmargin=1.4em, itemsep=2pt]
  \item Què li passa a l'ocupació a mesura que passa el temps?
  \item Cap a quin valor sembla que s'acosta? Arribarà mai a superar-lo?
  \item Com en diríeu, d'aquest valor cap al qual tendeix però no toca?
\end{itemize}
\end{idea}

\subsection{Exemples}

\begin{exemple}[la gràfica de l'ocupació]
La corba puja de pressa al principi i s'aplana a mesura que s'acosta a les $200$
places, sense traspassar-les mai:

\begin{center}
\begin{tikzpicture}
\begin{axis}[udaxis, width=10cm, height=5.4cm,
    xlabel={\small mesos}, ylabel={\small places ocupades},
    xmin=0, xmax=24, ymin=0, ymax=240,
    xtick={0,6,12,18,24}, ytick={0,100,200}]
  \addplot[udblue, domain=0:24, samples=100]{200*(1-exp(-x/6))};
  \addplot[udnavy, dashed, domain=0:24, samples=2]{200};
  \node[udnavy, font=\scriptsize, anchor=south east] at (axis cs:24,200) {capacitat $=200$};
\end{axis}
\end{tikzpicture}

\figcap{L'ocupació s'enfila cap a les $200$ places (el «sostre») però no les supera mai.}
\end{center}

\noindent La corba s'acosta cada cop més a la línia de les $200$ places, però no hi
arriba del tot: aquest valor és el \textbf{sostre} del servei.
\end{exemple}

\begin{exemple}[llegir valors concrets]
Si mirem l'ocupació mes a mes, veiem que els augments són cada cop més petits:
\begin{center}
\begin{tabular}{lcccccc}
\toprule
Mes & $3$ & $6$ & $12$ & $18$ & $24$ \\
\midrule
Places (aprox.) & $79$ & $126$ & $173$ & $190$ & $196$ \\
\bottomrule
\end{tabular}
\end{center}
Cada cop s'acosta més a $200$, però els salts es fan més petits: mai no la superarà.
\end{exemple}

\subsection{Exercicis resolts}

\begin{exresolt}[descriure el comportament]
Descriu amb les teves paraules què li passa a l'ocupació del centre a mesura que
passa el temps, i cap a quin valor tendeix.
\solucio
Al principi l'ocupació puja de pressa (s'omplen moltes places ràpidament). Després,
com que cada cop queden menys places lliures, l'augment es fa més lent i la corba
s'\textbf{aplana}. L'ocupació s'acosta cada cop més a les \textbf{$200$ places}, la
capacitat màxima, però no la supera mai. El valor $200$ és el \textbf{sostre} cap al
qual tendeix.
\resposta{Creix de pressa i després cada cop més a poc a poc, acostant-se a $200$
sense superar-les.}
\end{exresolt}

\begin{exresolt}[per què no se supera el sostre]
Per què té sentit, en aquest context real, que la gràfica no pugui passar de $200$?
\solucio
Perquè $200$ és la \textbf{capacitat física} del centre: no hi caben més de $200$
persones. Per molt que augmenti la demanda, l'ocupació no pot superar les places
disponibles. Per això la gràfica s'acosta a $200$ però mai la traspassa: el context
imposa aquest sostre.
\resposta{Perquè $200$ és la capacitat màxima del centre; no hi caben més places.}
\end{exresolt}

\subsection{Exercicis per fer}

\noindent\textit{Treballeu en grup a la pissarra, descrivint el comportament de cada
fenomen.}

\begin{exercicis}
\begin{exlist}
  \item Per a l'ocupació del centre, completa: «al principi creix \rule{2cm}{0.4pt} i
        després cada cop més \rule{2cm}{0.4pt}, acostant-se a \rule{1.5cm}{0.4pt} sense
        superar-les».
  \item Un altre centre té capacitat per a $150$ places. Si la seva ocupació es comporta
        igual (creix i s'aplana), cap a quin valor s'acostarà la gràfica?
  \item Mirant la taula de l'exemple, quants augments de places hi ha del mes $3$ al
        $6$, i del mes $18$ al $24$? Què observes en la mida d'aquests augments?
  \item \textbf{(Anticipació)} Si esperéssim molts mesos més, creus que l'ocupació
        arribaria \emph{exactament} a $200$? Justifica la teva intuïció.
  \item \textbf{(Vocabulari)} A les pissarres han sortit paraules per anomenar aquest
        valor cap al qual tendeix la gràfica («sostre», «topall», «màxim»\dots). Anota
        dues maneres diferents que hagin aparegut a classe.
  \item \textbf{(Context propi)} Pensa en un altre servei de l'àmbit social que
        s'acosti a un sostre (per exemple, una llista d'espera limitada o una sala amb
        aforament). Descriu-ne el comportament final.
\end{exlist}
\end{exercicis}

% ---------------------------------------------------------------------
\fullsolucions{Sessió 2}
\begin{solucions}
\begin{sollist}
  \item «al principi creix \textbf{de pressa} i després cada cop més \textbf{a poc a
        poc}, acostant-se a \textbf{$200$ places} sense superar-les».
  \item S'acostarà a $150$ places (la seva capacitat), sense superar-les.
  \item Del mes $3$ al $6$: $126-79=47$ places. Del mes $18$ al $24$: $196-190=6$
        places. Els augments es fan \textbf{cada cop més petits} a mesura que s'acosta
        al sostre.
  \item Resposta oberta. Idea esperada: s'hi acostaria moltíssim, gairebé fins a
        tocar-lo, però sense arribar-hi \emph{exactament}: sempre quedaria un marge
        diminut. (Aquesta és la idea de límit.)
  \item Resposta oberta (p.\,ex. «sostre», «topall», «capacitat», «màxim»).
  \item Resposta oberta.
\end{sollist}
\end{solucions}
